%
%  このファイルはヒューマンインタフェース研究報告集用クラスファイル
%  hisken.cls(ver.1.0) を利用した「原稿執筆の手引き」です．
%  Revised in Jun. 7, 2007 by Yutaka Ishii
%

\documentclass{hisken}
%\usepackage[dviout]{graphicx}
%\usepackage[dvipdfm]{graphicx}
\usepackage{graphicx}
\usepackage{epsfig} \let\epsfile=\epsfig


%和文タイトル
\jtitle{ヒューマンインタフェース研究報告集　原稿執筆の手引き}

%著者日本名
\jauthor{
ヒューマン　太郎\thanks{ヒューマン大学大学院　工学研究科}　　
ヒューマン　花子\thanks{ヒューマンインタフェース学会株式会社　中央研究所}　　
ヒューマン　次郎\addtocounter{footnote}{-2}\footnotemark
}

%英文タイトル
\etitle{Guideline of Writing Manuscript for Correspondences of Human Interface Society}

%著者英文名
\eauthor{
Taro Human\thanks{Graduate School of Engineering, Human University},
Hanako Human\thanks{Central Research Laboratory, Human Interface Society Ltd.} and 
Jiro Human\addtocounter{footnote}{-2}\footnotemark
}
% 名前は，はじめの一文字だけ大文字にしてください．

\begin{document}

%maketitle は abstract と keyword の後に入れてください．

\begin{abstract}
This paper describes the way how to write your manuscript for Correspondences of Human Interface Society.
\end{abstract}

\begin{keyword}	
keyword 1, keyword 2, keyword 3, keyword 4, keyword 5
\end{keyword}

\maketitle	

%%%%%%%%%%%%%%%%%%%%%%%%%%%%%%%%%%%%%%%%%%%%%%%%%%%%%%%%%%%%%%%%%%%%%
\section{はじめに}
%%%%%%%%%%%%%%%%%%%%%%%%%%%%%%%%%%%%%%%%%%%%%%%%%%%%%%%%%%%%%%%%%%%%%
ヒューマンインタフェース学会研究報告集へ掲載する原稿は，この「原稿執筆の手引き」に従って作成してください．なお，この「原稿執筆の手引き」自体が原稿のフォーマットに従っています．原稿執筆の際の参考にしてください．
また，本原稿に記載の研究開発は、学会が定める「ヒューマンインタフェース研究開発のための倫理指針」（https://www.his.gr.jp/office/ethical\_guidelines.html） を遵守して実施していることをご確認ください．とくに，原稿中に掲載される写真については，研究対象者の肖像権に配慮し，誰に対しどの場面を公開するかを研究対象者に示した上で諾否を得てください．


%%%%%%%%%%%%%%%%%%%%%%%%%%%%%%%%%%%%%%%%%%%%%%%%%%%%%%%%%%%%%%%%%%%%%
\section{題目など}
%%%%%%%%%%%%%%%%%%%%%%%%%%%%%%%%%%%%%%%%%%%%%%%%%%%%%%%%%%%%%%%%%%%%%

\subsection{題目}
和文原稿の場合は，和文題目と英文題目を併記してください．英文原稿の場合には英文題目のみを記載してください．

\subsection{著者}
和文原稿の場合は，和文著者名と英文著者名を併記してください．英文原稿の場合には英文著者名のみを記載してください．英文著者名は，名\,姓の順とし，それぞれ先頭の文字のみを大文字にしてください(ミドルネームがある場合は，名と姓の間にイニシャルを入れてください)．

\subsection{著者所属}
和文原稿の場合は，和文著者所属と英文著者所属を併記してください．英文原稿の場合には英文著者所属のみを記載してください．記載には，該当著者の右肩に$^ {*1,*2,\cdots}$などをつけ，第１頁左下に脚注としてそれぞれ記載してください．

\subsection{英文要旨}
論文の内容を簡潔にまとめた200語以内の英文要旨(Abstract)を記載してください．

\subsection{キーワード}
論文の内容を表すキーワードを5語以内で記載してください．


%%%%%%%%%%%%%%%%%%%%%%%%%%%%%%%%%%%%%%%%%%%%%%%%%%%%%%%%%%%%%%%%%%%%%
\section{本文}
%%%%%%%%%%%%%%%%%%%%%%%%%%%%%%%%%%%%%%%%%%%%%%%%%%%%%%%%%%%%%%%%%%%%%

\subsection{文章}
和文原稿の文章は口語体とし，特に英文もしくはカタカナ書きを必要とする部分以外は，現代かなづかいによる漢字かなまじり文としてください．外国の固有名詞は原則として原語綴りとします．
\subsection{用語}
学術用語は原則として文部省編の「学術用語集」，またはJISで制定されたものを用いてください．
\subsection{見出し}
本文の区分けはポイントシステムによる記号を用いて，章見出し「1., 2.,\,$ \cdots$\,」，節見出し「1.1, 1.2,\,$ \cdots$\,」，項見出し「1.1.1,1.1.2,\,$ \cdots$\,」などとしてください．付録は参考文献の後につけ，章見出しを無番号で「付録」とします．必要ならば節見出しとして「付録1.,付録2.,\,$ \cdots$\,」などを用いて区分けしてください．なお，章見出しは原稿の2行分を使用し，中央配置(センタリング)してください．

\subsection{文章の区切り}
文章の区切りにはピリオド「．」または句点「。」を，句の区切りにはコンマ「，」または読点「、」を用い，1字分をあててください．句読点は，どちらでも結構ですが，「ピリオド・コンマ」，あるいは「句点・読点」に統一してください．なお，各段落の最初は1字分をあけてください．


%%%%%%%%%%%%%%%%%%%%%%%%%%%%%%%%%%%%%%%%%%%%%%%%%%%%%%%%%%%%%%%%%%%%%
\section{数式など}
%%%%%%%%%%%%%%%%%%%%%%%%%%%%%%%%%%%%%%%%%%%%%%%%%%%%%%%%%%%%%%%%%%%%%

\subsection{数値・単位}
単位は原則として国際単位系(SI)を用い，数値はアラビア数字を使用してください．

\subsection{数式}
原則として，すべての数式に式番をつけてください．式番は通し番号とし，(1), (2), $ \cdots $ のように表します．参照する場合は(1)式, (2)式, $ \cdots $とします．分数式は，式として独立したものは
\begin{equation}
	\frac{a+b}{c+d}
	\label{formula}
\end{equation}
のように改行して書きますが，本文と同じ行の場合には $ (a+b) / (c+d) $のように書いてください．

\subsection{定理・定義・補題}
定義，定理，補題などの番号は通し番号とし，【定義1】，【定理1】，【補題1】$ \cdots $のように表します．参照する場合は括弧を取り，定義1，定理1などとします．


%%%%%%%%%%%%%%%%%%%%%%%%%%%%%%%%%%%%%%%%%%%%%%%%%%%%%%%%%%%%%%%%%%%%%
\section{図・表}
%%%%%%%%%%%%%%%%%%%%%%%%%%%%%%%%%%%%%%%%%%%%%%%%%%%%%%%%%%%%%%%%%%%%%

\subsection{番号}
図，表の番号は，それぞれ，図1,\,図2,\,$ \cdots $，Fig.1, \, Fig.2, \, $ \cdots $，表1,\,表2,\, $ \cdots $，Table\,1,\,Table\,2,\, $ \cdots $のように通し番号としてください．写真は図として扱います．

\subsection{見出し}
図の場合にはその下に，表の場合にはその上に，番号とともに見出しを入れてください．{\underline{和文原稿の場合に} \underline{は，日本語の見出し，英語の見出しの順で両方を入れ} \underline{てください．}} 英文原稿の場合には，英語の見出しのみを入れてください．

\subsection{引用}
本文中で図，表を引用する場合には，和文原稿の場合，それぞれ，図1,\,図2,\,$ \cdots $，表1,\,表2,\, $ \cdots $とし，英文原稿の場合は，それぞれ，Fig.1, \, Fig.2, \, $ \cdots $，Table\,1,\,Table\,2,\, $ \cdots $とします．

\subsection{サイズ}
図，表の刷り上がり寸法は，横幅8cm以内(片段)と横幅17cm以内(両段)の二通りとします．図や表中の文字は適切な大きさで，本文との整合性に注意してください．

\subsection{図，表の例}
図と表の記載例をこの「原稿執筆の手引き」中に示します（図\ref{fig:example1}，および表\ref{table:example2}参照)．

\begin{figure}[tb]
	\begin{center}
    \includegraphics[width=70mm]{fig1.eps}
%%    \includegraphics[width=70mm, bb=0 0 200 80]{fig1.bmp}
%%    \includegraphics[width=70mm, bb=0 0 200 80]{fig1.jpg} %%jpgを使用するためには、Susieプラグイン等が必要です
	\caption{楕円と五角形}		%和文 caption
	\ecaption{Ellipse and Pentagon.} %英文 caption
	\label{fig:example1}
	\end{center}
\end{figure}

\begin{table}[t]
	\begin{center}
	\caption{図形の辺と頂点}		%和文 caption
	\ecaption{Sides and Apices of Figures.} %英文 caption
	\label{table:example2}
	\begin{tabular}[hbt]{c c c}
	\hline
	\bf 図形 & \bf 辺の数 & \bf 頂点の数\\
	\hline
	三角形 & 3 & 3\\
	五角形 & 5 & 5\\
	楕円 & なし & なし\\
	\hline
	\end{tabular}
	\end{center}
\end{table}


%%%%%%%%%%%%%%%%%%%%%%%%%%%%%%%%%%%%%%%%%%%%%%%%%%%%%%%%%%%%%%%%%%%%%
\section{脚注}
%%%%%%%%%%%%%%%%%%%%%%%%%%%%%%%%%%%%%%%%%%%%%%%%%%%%%%%%%%%%%%%%%%%%%
脚注は本文の一部分として作成してください．

\subsection{引用}
脚注の引用は引用箇所の肩に $^ {1,2,3, \cdots} $ あるいは $^ {*,**,***,\cdots} $ ，$^ {\dag, \ddag, \S, \P} $などとつけてください\footnote{脚注はここに書いてください．}．

\subsection{脚注の記載}
脚注は，本文の下に境界を表す横線を引き，その下に記載してください．なお，脚注も原稿の枠内に書いてください．


%%%%%%%%%%%%%%%%%%%%%%%%%%%%%%%%%%%%%%%%%%%%%%%%%%%%%%%%%%%%%%%%%%%%%
\section{参考文献}
%%%%%%%%%%%%%%%%%%%%%%%%%%%%%%%%%%%%%%%%%%%%%%%%%%%%%%%%%%%%%%%%%%%%%
参考文献は本文の一部分として作成してください．

\subsection{引用}
参考文献の引用は引用箇所に$^ {\scriptsize [1],[2],[4]-[6]} $などとつけてください．

\subsection{文献の記載}
参考文献は本文の末尾にまとめてください．雑誌の場合は，著者名, 題目, 雑誌名(略記にて可), 巻(太字), 号, ページ, 発行西暦年を，書籍の場合には，著者(または編者)名, 書名(編者), 発行所, ページ,発行西暦年の順に記載してください．なお，著者(または編者)名の後はコロン「:」，題目の後はセミコロン「;」，そのほかはカンマ「,」で区切り，発行西暦年は小括弧「()」で囲んでください．

著者(または編者)名は，和文の場合，姓のみ羅列し，カンマ「,」で区切ってください．著者が多い場合には，代表的な著者名を記載し，その他の著者名を「他」で省略することができます．英文の場合には，「姓,名のイニシャル.」の羅列とし，カンマ「,」で区切ってください．著者が多い場合には，和文の場合と同様に「et al.」で省略することができます．\\

\noindent{（記載例）}
	{\small
	\begin{itemize}
	\item[{[1]}]
		西田, 井上, 吉川:
		ヒューマンインタフェース学会論文誌原稿執筆の手引き;
		ヒューマンインタフェース学会論文誌,
		{\bf Vol.1}, No.1, pp.1-10 (1999).
	\item[{[2]}]
		Nishida,S., Inoue,K., Yoshikawa,H.:
		Guideline of Writing Manuscript for the Transactions of Human Interface Society;
		The Transactions of Human Interface Society,
		{\bf Vol.1}, No.1, pp.1-10 (1999).
	\item[{[3]}]
		渋谷,高橋,他:
		21世紀の学会のあり方について;
		ヒューマンインタフェース学会誌,
		{\bf Vol.1}, No.1, pp.21-24 (1999).
	\item[{[4]}]
		高橋,渋谷,加藤:
		学会におけるネットワーク情報サービス;
		ヒューマンインタフェース(田村編),
		オーム社, 第25章 (1998).
	\item[{[5]}]
		Takahashi,M., Shibuya,Y., et al.:
		Network Information Service in Symposium;
		Human Interface (Tamura,H. ed.),
		Ohmsha, Chap.25 (1998).
	\end{itemize}}


%%%%%%%%%%%%%%%%%%%%%%%%%%%%%%%%%%%%%%%%%%%%%%%%%%%%%%%%%%%%%%%%%%%%%
\section{文字数・大きさ}
%%%%%%%%%%%%%%%%%%%%%%%%%%%%%%%%%%%%%%%%%%%%%%%%%%%%%%%%%%%%%%%%%%%%%

\subsection{文字数}
原稿は，A4サイズの研究報告集に印刷されます．
本文の文字数は，25文字×49行の2段組とします．各部のマージン(余白)は表\ref{table:margin}の通りです．

\begin{table}[bt]
	\begin{center}
	\caption{原稿のマージン}		%和文 caption
	\ecaption{Margins of Manuscript.} 	%英文 caption
	\label{table:margin}
	\begin{tabular}[hbt]{c r}
	\hline
	\bf 各部 & \bf マージン\\
        \hline
        用紙上端からタイトル(1頁目) & 37 mm\\
        用紙上端から本文(2頁目以降) & 30 mm\\
        用紙下端から本文(全頁) & 15 mm\\
        用紙左端から本文(全頁) & 20 mm\\
        用紙右端から本文(全頁) & 20 mm\\
        段組中央(全頁) & 10 mm\\
        \hline
	\end{tabular}
	\end{center}
\end{table}

\subsection{文字の大きさ}
文字の大きさは，原則として表\ref{table:charsize}の通りとします．

\begin{table}[tbh]
	\begin{center}
	\caption{文字の大きさ}		%和文 caption
	\ecaption{Character Size.} 	%英文 caption
	\label{table:charsize}
	\begin{tabular}[hbt]{c c r}
	\hline
	\bf 文字 & \bf 書体 & \bf ポイント\\
	\hline
	和文表題 & ゴチック体 & 16.5\\
	和文著者名 & 明朝体 & 14.0\\
	英文表題(和文原稿) & Times(Bold) & 12.0\\
	英文著者名(和文原稿) & Times & 12.0\\
	英文表題(英文原稿) & Times(Bold) & 14.0\\
	英文著者名(英文原稿) & Times & 12.0\\
	``Abstract'' & Times(Bold) & 9.0\\
	英文要旨本文 & Times & 9.0\\
	``Keywords'' & Times(Bold) & 9.0 \\
	キーワード & Times & 9.0\\ \vspace{-0.7mm}
	章・節見出し & ゴチック体 & 9.0\\
	& Times(Bold) & \\
	項見出し & 明朝体・Times & 9.0\\
	図等見出し & 明朝体・Times & 9.0\\
	本文 & 明朝体・Times & 9.0\\
	参考文献 & 明朝体・Times & 8.0\\ \vspace{-0.7mm}
	脚注 & 明朝体・Times &  8.0\\
	\hline
	\end{tabular}
	\end{center}
\end{table}



%%%%%%%%%%%%%%%%%%%%%%%%%%%%%%%%%%%%%%%%%%%%%%%%%%%%%%%%%%%%%%%%%%%%%
\section{原稿の作成}
%%%%%%%%%%%%%%%%%%%%%%%%%%%%%%%%%%%%%%%%%%%%%%%%%%%%%%%%%%%%%%%%%%%%%

ヒューマンインタフェース学会のWeb Siteから，「原稿執筆の手引き」に準拠したフォーマットのMicrosoft Wordで作成されたファイル（meeting.doc）がダウンロードできます．

\LaTeX2.09{}\cite{tex1,tex2}により原稿を作成する場合には，スタイルファイル{\tt hisken.sty} (和文原稿用)，または{\tt ehisken.sty} (英文原稿用)を利用してください．また，\LaTeX{}2$\varepsilon$\cite{tex3}により原稿を作成する場合には，クラスファイル{\tt hisken.cls} (和文原稿用)，または{\tt ehisken.cls} (英文原稿用)を利用してください．これらをドキュメントスタイルとして指定することで，この「原稿執筆の手引き」に準拠したフォーマットの原稿を簡単に作成することができます．スタイルファイルは，ヒューマンインタフェース学会のWeb Siteからダウンロードすることができます．

なお，作成されました原稿はAcrobat PDF（非圧縮フォント埋め込み）へ変換し，必ず各自で体裁等ご確認の上，3.5インチフロッピーディスク，CDメディア，DVDメディアのいずれかにディレクトリ階層なしで記録して学会事務局に郵送頂くか，メールの添付ファイルとして提出してください．


%%%%%%%%%%%%%%%%%%%%%%%%%%%%%%%%%%%%%%%%%%%%%%%%%%%%%%%%%%%%%%%%%%%%%
\section*{謝辞}
%%%%%%%%%%%%%%%%%%%%%%%%%%%%%%%%%%%%%%%%%%%%%%%%%%%%%%%%%%%%%%%%%%%%%
この「原稿執筆の手引き」を作成するにあたり，日本計測自動制御学会をはじめとする諸先輩学会の原稿執筆の手引きを参考にさせていただきました．ここに感謝の意を表します．また，\LaTeX{}のスタイルファイルを作成するにあたり，その元となるスタイルファイルを快く提供していただいた日本VR学会に深く感謝いたします．

%%%%%%%%%%%%%%%%%%%%%%%%%%%%%%%%%%%%%%%%%%%%%%%%%%%%%%%%%%%%%%%%%%%%%
% 　参考文献
%%%%%%%%%%%%%%%%%%%%%%%%%%%%%%%%%%%%%%%%%%%%%%%%%%%%%%%%%%%%%%%%%%%%%
\begin{thebibliography}{99}
\bibitem{tex1}
	Knuth,D.E.:
	\TeX{}ブック,
	アスキー出版局,(1989).
\bibitem{tex2}
	奥村:
	\LaTeX{}美文書作成入門,
	技術評論社,(1991).
\bibitem{tex3}
    奥村:
    \LaTeX2$\varepsilon$美文書作成入門（改訂第4版）,
    技術評論社,(2006).
\end{thebibliography}


%%%%%%%%%%%%%%%%%%%%%%%%%%%%%%%%%%%%%%%%%%%%%%%%%%%%%%%%%%%%%%%%%%%%%
\appendix{}
%%%%%%%%%%%%%%%%%%%%%%%%%%%%%%%%%%%%%%%%%%%%%%%%%%%%%%%%%%%%%%%%%%%%%
付録は参考文献の後につけ，章見出しを無番号で「付録」とします．必要ならば節見出しとして「付録1.,付録2.,\,$ \cdots $\,」などを用いて区分けしてください．

\vspace{15mm}

\section*{改定履歴}
\subsection*{Ver.1.1 on 1999/July/7}
\begin{itemize}
\item 原稿のヘッダー情報を削除
\item 英文著者名にミドルネームの記載を追加
\item 原稿フォーマットテンプレートファイルの記載を追加
\item 第１頁の上端マージンを変更
\end{itemize}

\subsection*{Ver.1.2 on 2007/June/13}
\begin{itemize}
\item 原稿作成方法，提出方法の修正
\item \LaTeX2$\varepsilon$に関する説明の追加
\item 参考文献の追加
\end{itemize}

\subsection*{Ver.1.3 on 2015/July/28}
\begin{itemize}
\item 倫理指針遵守に関する説明の追加
\end{itemize}


\end{document}

